% Options for packages loaded elsewhere
\PassOptionsToPackage{unicode}{hyperref}
\PassOptionsToPackage{hyphens}{url}
\PassOptionsToPackage{dvipsnames,svgnames,x11names}{xcolor}
\documentclass[
  10pt,
]{article}
\usepackage{xcolor}
\usepackage[margin=2.2cm]{geometry}
\usepackage{amsmath,amssymb}
\setcounter{secnumdepth}{-\maxdimen} % remove section numbering
\usepackage{iftex}
\ifPDFTeX
  \usepackage[T1]{fontenc}
  \usepackage[utf8]{inputenc}
  \usepackage{textcomp} % provide euro and other symbols
\else % if luatex or xetex
  \usepackage{unicode-math} % this also loads fontspec
  \defaultfontfeatures{Scale=MatchLowercase}
  \defaultfontfeatures[\rmfamily]{Ligatures=TeX,Scale=1}
\fi
\usepackage{lmodern}
\ifPDFTeX\else
  % xetex/luatex font selection
\fi
% Use upquote if available, for straight quotes in verbatim environments
\IfFileExists{upquote.sty}{\usepackage{upquote}}{}
\IfFileExists{microtype.sty}{% use microtype if available
  \usepackage[]{microtype}
  \UseMicrotypeSet[protrusion]{basicmath} % disable protrusion for tt fonts
}{}
\makeatletter
\@ifundefined{KOMAClassName}{% if non-KOMA class
  \IfFileExists{parskip.sty}{%
    \usepackage{parskip}
  }{% else
    \setlength{\parindent}{0pt}
    \setlength{\parskip}{6pt plus 2pt minus 1pt}}
}{% if KOMA class
  \KOMAoptions{parskip=half}}
\makeatother
\usepackage{color}
\usepackage{fancyvrb}
\newcommand{\VerbBar}{|}
\newcommand{\VERB}{\Verb[commandchars=\\\{\}]}
\DefineVerbatimEnvironment{Highlighting}{Verbatim}{commandchars=\\\{\}}
% Add ',fontsize=\small' for more characters per line
\newenvironment{Shaded}{}{}
\newcommand{\AlertTok}[1]{\textcolor[rgb]{1.00,0.00,0.00}{\textbf{#1}}}
\newcommand{\AnnotationTok}[1]{\textcolor[rgb]{0.38,0.63,0.69}{\textbf{\textit{#1}}}}
\newcommand{\AttributeTok}[1]{\textcolor[rgb]{0.49,0.56,0.16}{#1}}
\newcommand{\BaseNTok}[1]{\textcolor[rgb]{0.25,0.63,0.44}{#1}}
\newcommand{\BuiltInTok}[1]{\textcolor[rgb]{0.00,0.50,0.00}{#1}}
\newcommand{\CharTok}[1]{\textcolor[rgb]{0.25,0.44,0.63}{#1}}
\newcommand{\CommentTok}[1]{\textcolor[rgb]{0.38,0.63,0.69}{\textit{#1}}}
\newcommand{\CommentVarTok}[1]{\textcolor[rgb]{0.38,0.63,0.69}{\textbf{\textit{#1}}}}
\newcommand{\ConstantTok}[1]{\textcolor[rgb]{0.53,0.00,0.00}{#1}}
\newcommand{\ControlFlowTok}[1]{\textcolor[rgb]{0.00,0.44,0.13}{\textbf{#1}}}
\newcommand{\DataTypeTok}[1]{\textcolor[rgb]{0.56,0.13,0.00}{#1}}
\newcommand{\DecValTok}[1]{\textcolor[rgb]{0.25,0.63,0.44}{#1}}
\newcommand{\DocumentationTok}[1]{\textcolor[rgb]{0.73,0.13,0.13}{\textit{#1}}}
\newcommand{\ErrorTok}[1]{\textcolor[rgb]{1.00,0.00,0.00}{\textbf{#1}}}
\newcommand{\ExtensionTok}[1]{#1}
\newcommand{\FloatTok}[1]{\textcolor[rgb]{0.25,0.63,0.44}{#1}}
\newcommand{\FunctionTok}[1]{\textcolor[rgb]{0.02,0.16,0.49}{#1}}
\newcommand{\ImportTok}[1]{\textcolor[rgb]{0.00,0.50,0.00}{\textbf{#1}}}
\newcommand{\InformationTok}[1]{\textcolor[rgb]{0.38,0.63,0.69}{\textbf{\textit{#1}}}}
\newcommand{\KeywordTok}[1]{\textcolor[rgb]{0.00,0.44,0.13}{\textbf{#1}}}
\newcommand{\NormalTok}[1]{#1}
\newcommand{\OperatorTok}[1]{\textcolor[rgb]{0.40,0.40,0.40}{#1}}
\newcommand{\OtherTok}[1]{\textcolor[rgb]{0.00,0.44,0.13}{#1}}
\newcommand{\PreprocessorTok}[1]{\textcolor[rgb]{0.74,0.48,0.00}{#1}}
\newcommand{\RegionMarkerTok}[1]{#1}
\newcommand{\SpecialCharTok}[1]{\textcolor[rgb]{0.25,0.44,0.63}{#1}}
\newcommand{\SpecialStringTok}[1]{\textcolor[rgb]{0.73,0.40,0.53}{#1}}
\newcommand{\StringTok}[1]{\textcolor[rgb]{0.25,0.44,0.63}{#1}}
\newcommand{\VariableTok}[1]{\textcolor[rgb]{0.10,0.09,0.49}{#1}}
\newcommand{\VerbatimStringTok}[1]{\textcolor[rgb]{0.25,0.44,0.63}{#1}}
\newcommand{\WarningTok}[1]{\textcolor[rgb]{0.38,0.63,0.69}{\textbf{\textit{#1}}}}
\usepackage{longtable,booktabs,array}
\usepackage{caption}
\captionsetup[table]{skip=6pt}
\newcounter{none} % for unnumbered tables
\usepackage{calc} % for calculating minipage widths
% Correct order of tables after \paragraph or \subparagraph
\usepackage{etoolbox}
\makeatletter
\patchcmd\longtable{\par}{\if@noskipsec\mbox{}\fi\par}{}{}
\makeatother
\usepackage{graphicx}
\makeatletter
\newsavebox\pandoc@box
\newcommand*\pandocbounded[1]{% scales image to fit in text height/width
  \sbox\pandoc@box{#1}%
  \Gscale@div\@tempa{\textheight}{\dimexpr\ht\pandoc@box+\dp\pandoc@box\relax}%
  \Gscale@div\@tempb{\linewidth}{\wd\pandoc@box}%
  \ifdim\@tempb\p@<\@tempa\p@\let\@tempa\@tempb\fi% select the smaller of both
  \ifdim\@tempa\p@<\p@\scalebox{\@tempa}{\usebox\pandoc@box}%
  \else\usebox{\pandoc@box}%
  \fi%
}
% Set default figure placement to htbp
\def\fps@figure{htbp}
\makeatother
\setlength{\emergencystretch}{3em} % prevent overfull lines
\providecommand{\tightlist}{%
  \setlength{\itemsep}{0pt}\setlength{\parskip}{0pt}}
\usepackage{fontspec}
\setmainfont{STIX Two Text}
\setmonofont{Menlo}[Scale=0.85]
\usepackage{xeCJK}
\setCJKmainfont{Songti SC}
\usepackage{microtype}
\usepackage{etoolbox}
\AtBeginEnvironment{longtable}{\footnotesize}
\usepackage{xurl}
\setlength{\emergencystretch}{3em}
\sloppy
\setCJKmonofont{Songti SC}
\usepackage{unicode-math}
\setmathfont{STIX Two Math}
\usepackage{newunicodechar}
\newunicodechar{→}{\ensuremath{\rightarrow}}
\newunicodechar{∎}{\ensuremath{\blacksquare}}
\newunicodechar{≠}{\ensuremath{\neq}}
\newunicodechar{≈}{\ensuremath{\approx}}
\newunicodechar{↔}{\ensuremath{\leftrightarrow}}
\newunicodechar{≪}{\ensuremath{\ll}}
\newunicodechar{↛}{\ensuremath{\nrightarrow}}
\newunicodechar{・}{\textperiodcentered}
\newunicodechar{⇒}{\ensuremath{\Rightarrow}}
\newunicodechar{∪}{\ensuremath{\cup}}
\newunicodechar{≥}{\ensuremath{\geq}}
\newunicodechar{≤}{\ensuremath{\leq}}
\newunicodechar{⟨}{\ensuremath{\langle}}
\newunicodechar{⟩}{\ensuremath{\rangle}}
\usepackage{bookmark}
\IfFileExists{xurl.sty}{\usepackage{xurl}}{} % add URL line breaks if available
\urlstyle{same}
% fallback for those not using the hyperref driver hyperxmp:
\makeatletter
\@ifundefined{xmpquote}{\newcommand{\xmpquote}[1]{#1}}{}
\makeatother
\hypersetup{
  pdftitle={Task Satiation vs.~Aggregate Rebound in AI Demand},
  pdfauthor={David Charlot, Open Interface Engineering},
  colorlinks=true,
  linkcolor={blue},
  filecolor={Maroon},
  citecolor={Blue},
  urlcolor={blue},
  pdfcreator={LaTeX via pandoc}}

\title{Task Satiation vs.~Aggregate Rebound in AI Demand}
\author{David Charlot, Open Interface Engineering}
\date{Version 1.1, 11 October 2026. HTML:
https://research.openie.dev/papers/satiation/}

\begin{document}
\maketitle

{
\setcounter{tocdepth}{2}
\tableofcontents
}
\emph{Version 1.1, 11 October 2026. Data and code:
\href{https://research.openie.dev/data/\#economics}{research.openie.dev/data/}.
LaTeX PDF:
\href{https://research.openie.dev/pdfs/latex/satiation.pdf}{/pdfs/latex/satiation.pdf}
(\href{https://research.openie.dev/pdfs/latex/satiation.tex}{source}).}

\subsection{Abstract}\label{abstract}

Once a task is done, more machine effort on that task adds nothing. This
paper defines that stop precisely and separates it from the separate
question of whether total demand for machine work keeps growing.

\textbf{In plain words}

\begin{enumerate}
\def\labelenumi{\arabic{enumi}.}
\tightlist
\item
  Every job has a finish line. After it, more computer work on that job
  is waste.
\item
  Cheaper computer work can still make people start many more jobs, so
  total use can keep rising.
\item
  Money prices for AI are falling fast, but the electricity it uses is
  never free.
\end{enumerate}

Many planning models treat demand for machine intelligence as if it had
no finish line. Consumer theory already contains the opposite object.
Satiation, and the related bliss point, is the region in which further
units of a good do not raise utility and may lower it (Andersen, 2001).
This paper defines an operational cousin for work and care: a loop
reaches task completion when a stated completeness predicate holds,
further calls do not flip any clause of that predicate from false to
true, and calls that undo the predicate are harm, not value. Engagement
time, token count, and seat count are different variables. They can rise
after completeness.

Two claims must be kept apart. Task-level satiation holds by definition:
once a task's completeness predicate is true, more synthesis on that
task buys nothing. Aggregate satiation is a claim about total demand,
and it is not implied. Cheaper synthesis can open new tasks (Acemoglu
and Restrepo, 2019) and raise total use, as cheaper coal raised total
coal use (Jevons, 1865). Each chore can be satiated while total demand
for chores grows without bound. This rebound is the central
identification problem for the paper. Section 5.1 states two testable
hypotheses and the data that would identify them.

A second fact is about price, not physics. Epoch AI reports large
declines in the price of a fixed inference performance (Epoch AI, 2025;
Emberson and Roodman, 2026). This paper calls the resulting situation
free at the margin for digital chores whose output is information and
whose completeness test is finite. The phrase does not mean that joules
are free, that a plant may move without a certificate, or that care
labor has been automated away. Landauer's bound is a lower bound on
erasure, not a description of current accelerators (Landauer, 1961;
Horowitz, 2014).

The paper does not report a randomized trial, a done-detector error
rate, or a split between energy spent before completeness and energy
spent after it. Those measurements are specified in Section 3.3 and are
absent.

\subsection{Notation}\label{notation}

{\def\LTcaptype{none} % do not increment counter
\begin{longtable}[]{@{}
  >{\raggedright\arraybackslash}p{(\linewidth - 2\tabcolsep) * \real{0.1875}}
  >{\raggedright\arraybackslash}p{(\linewidth - 2\tabcolsep) * \real{0.8125}}@{}}
\toprule\noalign{}
\begin{minipage}[b]{\linewidth}\raggedright
Term
\end{minipage} & \begin{minipage}[b]{\linewidth}\raggedright
Definition in this paper
\end{minipage} \\
\midrule\noalign{}
\endhead
\bottomrule\noalign{}
\endlastfoot
Satiation & In the cited theory, a fall of marginal utility to zero or
below. In this paper's operational clause, after completeness, further
calls do not flip any clause of \texttt{C} from false to true; calls
that undo \texttt{C} are harm, not value (Section 3.1). \\
Bliss point & A bundle at which additional quantity of the good reduces
utility. A textbook object. Not estimated here. \\
Task completion & The event that the completeness predicate becomes
true. Orthogonal to model scale. \\
Work completeness & A stated work stop: ticket closed, patch merged,
report filed, ledger reconciled, or an equivalent written test. \\
Care completeness & A stated care stop: assessment finished, medications
reconciled, receiving party notified, person in a safe state, or an
equivalent written test. \\
Free at the margin & Marginal price of a digital synthesis step trends
toward bundled or negligible money price. Joules and commits stay
costly. \\
Artificial general intelligence (AGI) & A broad capability target. Not a
price observation and not a result of this paper. \\
Direct rebound & Extra use of a service when its effective price falls,
on the same service (Gillingham, Rapson, and Wagner). \\
Modeled / simulated / board-measured & Formula, program execution, or
meter. This paper's new numbers are none of the three; its empirical
content is citation of published series and documents. \\
\end{longtable}
}

Scope. Section 2 reviews satiation theory, task automation, attention
markets, technology-cycle evidence, and inference prices. Section 3
writes predicates and says which were not estimated. Section 4 lists
claims that rest on citations, and separates them from design rules.
Section 5 answers rebound, care costs, and military demand. Section 6
lists threats to validity.

\subsection{1. Introduction}\label{introduction}

A demand curve with constant positive slope in ``intelligence'' will
recommend more seats, more tokens, and more agents after a job is
finished. That recommendation is a modeling choice. It is not a theorem
about work. Standard consumer theory allows marginal utility to decline
and to cross zero. If a chore has a completeness predicate, the relevant
question is whether more model output changes that predicate. If it does
not, extra output is a cost with no increment on the objective that was
written down.

The recent price evidence points the other way from scarcity of
synthesis. Epoch's series show the money price of a given benchmark
performance falling quickly. Open weights and typed decision procedures,
including System One, are additional routes by which a previously scarce
generation step becomes cheap. Teams that price a closed chore as if
synthesis were still scarce are using a stale constraint.

The constraints that do not follow from a price series are physical and
institutional. Computation dissipates energy. Horowitz (2014) documents
that data movement and memory dominate practical CMOS energy, at scales
far above Landauer's ideal erasure bound. An actuator is not a token. A
signature, a clinical handoff, and a payment are not reversed by
sampling another paragraph. Baumol (1967) showed how sectors with weak
productivity growth can absorb a rising share of expenditure. Acemoglu
and Restrepo (2019) model automation as task displacement plus the
possible creation of new tasks. None of these results says that a
language model ends cost disease. They do say that ``more capability''
is not a sufficient description of value.

This paper's positive claim is definitional and bibliographic. Satiation
is a coherent object. Published prices for inference have fallen
sharply. Engagement metrics are not completeness metrics. A stop for
``do not spend'' and a stop for ``do not move'' are different
predicates. The negative claim is equally important. This paper does not
identify a structural demand system, does not estimate rebound, and does
not report measured energy.

\subsection{2. Related work}\label{related-work}

\subsubsection{2.1 Satiation and bliss
points}\label{satiation-and-bliss-points}

Andersen (2001) treats satiation inside an evolutionary model of
structural change, not as a verbal preference. The durable element for
this paper is the economic one: marginal valuation of a good can reach
zero. Textbook bliss points make the geometry explicit. For a bundle
\texttt{x*}, utility does not increase in a coordinate past \texttt{x*},
and a smooth representation can have negative marginal utility there. A
quadratic illustration, \texttt{U(x)\ =\ -a\ (x\ -\ x*)\^{}2} for
\texttt{a\ \textgreater{}\ 0}, has a maximum at \texttt{x*} and a
negative derivative beyond it. That function is a model used to teach
the definition. It is not fitted to agent logs in this study. No
coefficient \texttt{a} is reported.

The operational translation requires a predicate the model does not know
unless someone states it. Let \texttt{C(z)\ =\ 1} when the written
completeness test for episode \texttt{z} holds, else 0. Satiation on
that episode means there exists a time \texttt{t*} such that \texttt{C}
is 1 at \texttt{t*} and, for all later calls, no clause of \texttt{C}
flips from false to true. Calls that undo \texttt{C} are harm, not
value. Section 3.1 states the same rule. The internal event
\textbf{economic done} is when \texttt{C} becomes 1; this paper calls
that event task completion. This is a definition. It becomes an
empirical claim only with a labeled set of episodes and a detector whose
errors are counted. Section 4 does not contain that count.

Keynes drew the same line in 1930. He split needs into absolute needs,
felt whatever others have, and relative needs, felt because satisfying
them lifts us above others. He argued the absolute needs could be met.
The completeness predicate is the absolute need of one task, written
down. Relative and open-ended wants have no such predicate, and Section
5.2 keeps them out of the claim.

\subsubsection{2.2 Tasks, displacement, and cost
disease}\label{tasks-displacement-and-cost-disease}

Baumol (1967) analyzes unbalanced growth. If one sector's productivity
rises and another's does not, the relative cost of the stagnant sector
rises, and it can dominate expenditure even when its quantity grows
slowly. Health care, education, and live performance are the usual
examples. Later health-economics papers dispute magnitudes and
mechanisms. This paper uses Baumol for a negative instruction: do not
infer, from a fall in token price, that care has become a software good.
Care remains dependent on time, trust, and liability in the ordinary
institutional sense. A documentation model can increase recorded
activity, which is a Baumol-relevant ambiguity: measured output can move
while the care predicate does not.

The labor literature associated with David Autor decomposes jobs into
tasks that machines may or may not perform. The pinned statement is
Acemoglu and Restrepo (2019), ``Automation and New Tasks: How Technology
Displaces and Reinstates Labor,'' Journal of Economic Perspectives
33(2), https://www.aeaweb.org/articles?id=10.1257/jep.33.2.3. They argue
that displacement reduces labor's task share and that new tasks can
restore it. A completeness-oriented stop is closer to finishing a task
than to creating a new engagement task. An agent loop with no stop can
raise activity measures without a productivity gain. That possibility is
consistent with their warning about weak automation.

\subsubsection{2.3 Attention and
engagement}\label{attention-and-engagement}

Wu's book ``The Attention Merchants'' (publisher page:
https://www.penguinrandomhouse.com/books/234876/the-attention-merchants-by-tim-wu/)
describes businesses that harvest attention and resell it. The Center
for Humane Technology's overview of the attention economy
(https://www.humanetech.com/youth/the-attention-economy) is an advocacy
summary of the same pattern, not a trial. The inference used here is
limited. If a product's objective is time on task, then a session that
continues after \texttt{C\ =\ 1} can look successful. If the objective
is \texttt{C}, the same session is waste or harm. The two objectives are
not observationally equivalent in a token log. No experiment in this
paper estimates the welfare gap.

\subsubsection{2.4 Historical cycles}\label{historical-cycles}

Technology markets often show a speculative phase, a use phase, and a
later commodity phase. The evidence is uneven.

Quinn (2019), ``Technological revolutions and speculative finance:
evidence from the British Bicycle Mania,'' Cambridge Journal of
Economics 43(2), 271-294, https://doi.org/10.1093/cje/bey029, studies
cycle-share prices from 1895 to 1900. Prices rose by over 200 percent
and then fell by more than 75 percent. The paper's result is about
speculative finance and fundamentals, not about a measured ``enough
bicycle'' predicate for households. Using it as a picture of race-phase
finance is supported. Using it as a proof of transport satiation is not.

Cars, scheduled aviation, and television have commodity-like markets in
the everyday sense: used vehicles, yield-managed seats, and panels sold
as inputs to content. The rows are historical orientation. They are not
regressions. Airline yield management after deregulation is a pricing
institution for a seat, which is evidence of a commodity market in
seats, not evidence that passengers' demand for travel is globally
satiated.

Pew Research Center's mobile fact sheet
(https://www.pewresearch.org/internet/fact-sheet/mobile/) is the
saturation source used for phones. The project archive records U.S.
smartphone ownership near 91 percent in the 2024 and 2025 fact-sheet
readings. A reader who quotes the figure should open the live sheet,
because fact sheets are revised. High ownership is evidence of
diffusion. It is not, by itself, evidence that marginal value of another
phone camera is zero.

The digital-chore row is the one this paper defends with a price series
rather than with an analogy. See Section 4.2.

\subsubsection{2.5 Inference prices and open
weights}\label{inference-prices-and-open-weights}

Epoch AI (2025), ``LLM inference price trends,''
https://epoch.ai/data-insights/llm-inference-price-trends, reports price
declines at fixed performance on the order of 9 to 900 times per year
across benchmarks in that note. Emberson and Roodman for Epoch AI
(2026), ``The plunging price of thought,''
https://epoch.ai/publications/the-plunging-price-of-thought, report
about a 47 percent decline per quarter in the cost of a given
performance since about 2023, summarized as about 13 times per year,
faster near the frontier. These are the authors' summaries of their
data. This paper does not re-estimate the index. The economic reading is
narrow: the money price of synthesis for a fixed measured performance
has been falling fast. A falling supply price shifts quantity demanded
if demand is elastic. It does not create a new completeness predicate,
and it does not set the energy price to zero.

System One-class tools are relevant as a mechanism for a further price
cut on decision problems with a known option set. A System One model
returns a typed answer and a probability for each allowed option in one
forward pass, with no free-form text. TypeSafe AI named the category
with its hosted Jev model (15 September 2026). Open models followed,
among them Laya from Convai Innovations. A System One model proposes. It
is not a completeness evaluation.

A planning rule of thumb that frontier competence reaches cheaper
devices in about six months is not a measured constant in any source
cited here. Distillation, quantization, and table-lookup inference (Ma
et al., 2024; Wei et al., 2024) are published mechanisms by which a
given competence can move to cheaper hardware. The speed is an empirical
question per model pair.

\subsubsection{2.6 Physical cost}\label{physical-cost}

Landauer (1961) bounds erasure. Bennett (1973) shows reversible
computation can avoid that bound in theory. Horowitz (2014) is the
practical CMOS statement: energy per operation and energy of data
movement set the design constraint, at many orders of magnitude above
\texttt{kT\ ln\ 2} for irreversible CMOS logic. At \(T = 300\) K the
bound is \(k_B T \ln 2 \approx 2.87 \times 10^{-21}\) J per erased bit,
with \(k_B = 1.380649 \times 10^{-23}\) J/K.

Goldman Sachs published a research note titled ``Gen AI: Too Much Spend,
Too Little Benefit?''
(https://www.goldmansachs.com/insights/top-of-mind/gen-ai-too-much-spend-too-little-benefit).
The existence of a major sell-side note that questions infrastructure
spend relative to benefit is a document fact. Dollar figures inside the
note are not quoted here, because sell-side estimates are not a
scientific energy account.

\subsection{3. Definitions and methods}\label{definitions-and-methods}

\subsubsection{3.1 Predicates}\label{predicates}

Let an episode be a sequence of model or tool calls
\texttt{z\_1,\ ...,\ z\_T} together with a human-facing state. A
completeness predicate \texttt{C} maps the episode prefix to
\texttt{\{0,\ 1\}}. Examples that can be written without a model score:

\begin{itemize}
\tightlist
\item
  Work: tests pass and the change is merged; or the ledger discrepancy
  is zero; or the required fields of a report are filled and the report
  is filed.
\item
  Care: the medication list matches the source, the receiving clinician
  has the note, and the safety checklist items that the institution
  marked critical are true.
\end{itemize}

\texttt{C} is not a probability. A detector that outputs a score still
needs a threshold and an error table before it is a predicate in the
sense of Section 4's missing experiment.

The internal event \textbf{economic done} at time \texttt{t*} means
\texttt{C\ =\ 0} on the prefix before \texttt{t*} and \texttt{C\ =\ 1}
at \texttt{t*}, for the chosen test. This paper calls the catalog layer
\textbf{task completion}. Satiation relative to model calls means that
for all \texttt{t\ \textgreater{}\ t*}, extra calls do not flip any
clause of \texttt{C} from false to true. Calls that undo \texttt{C} are
harm, not value. The definition does not say humans always know
\texttt{t*}. It says that if no \texttt{C} is written, the satiation
claim is not yet a statement about that product.

Free at the margin is a statement about money price. Let
\texttt{p\_t(q)} be the price of synthesis service \texttt{q} at date
\texttt{t}. The archival price evidence is that \texttt{p\_t} for fixed
benchmark performance has fallen at the rates Epoch reports. Free at the
margin means \texttt{p\_t} is low enough that a seat license is a weak
explanation of cost for that chore. It does not mean the integral of
power over time is zero.

A stop rule can therefore have two kinds of reason, economic (the
predicate holds, or a budget or policy limit is reached) and physical (a
safety or energy limit is reached). The two kinds are different
predicates and should be recorded as different reasons.

\subsubsection{3.2 What was done in this
study}\label{what-was-done-in-this-study}

This study is a structured reading. Sources are the citation spine
checked in the project on 29 to 30 September 2026, and the satiation
notes. Section 3.3 names the measurements that would make the
operational definition empirical.

Design rules in Section 4.10 are recommendations conditional on the
definitions. They are not estimated treatment effects. A reader can
reject a rule without rejecting Epoch's price series, and can accept the
price series without accepting a product recommendation.

\subsubsection{3.3 Measurements specified and not
taken}\label{measurements-specified-and-not-taken}

\textbf{The predicate is endogenous.} Witt (2001) argues that wants are
learned, so the completeness predicate \(C(z)\) can ratchet as price
falls. A study of satiation across a price change must therefore freeze
\(C(z)\) in its preregistration before the change and report any change
of predicate as a separate treatment.

\textbf{H1 can be read from public data now.} Chen and colleagues (2024)
measure tokens spent after the first correct answer; Cuadron and
colleagues (2025) score overthinking in agent trajectories; Bai, Huang
and Wang (2026) find that same-task token use in agentic coding varies
up to 30 times across runs. Applying this paper's \(C(z)\) post hoc to
those trajectories gives an exploratory H1 estimate, labeled
exploratory. It is the next measurement.

\textbf{Tokens and joules.} H1 and H2 are stated in tokens; the thesis
is in joules. Masanet and colleagues (2020) found that data-centre
compute grew about sixfold from 2010 to 2018 while energy rose far less,
so a token elasticity is not an energy elasticity. Each hypothesis is to
be reported in both units, with the unit the data identify named.

\begin{enumerate}
\def\labelenumi{\arabic{enumi}.}
\tightlist
\item
  Split account. For each episode, sum analytical or, later,
  board-measured energy on calls with \texttt{t\ \textless{}=\ t*} and
  on calls with \texttt{t\ \textgreater{}\ t*}. Direct rebound on that
  chore is visible only if the post-done integral is estimated. The
  repository does not contain this split.
\item
  Detector errors. On a labeled corpus, false done (predicate declared
  true when a clause is false) and false continue (\texttt{C} true and
  the system keeps calling). No rates are published here.
\item
  Care handoff. Time to safe state and a handoff checklist, contrasted
  with session length. Not measured.
\item
  Segment table. Which products have a short \texttt{C} and which are
  open-ended research or adversarial races. Not estimated. The military
  case is excluded from the satiation claim in Section 5 rather than
  forced into it.
\item
  Measured energy. An external power meter on a stated workload. Not
  done.
\end{enumerate}

Until those exist, sentences about ``overkill joules'' are hypotheses.
Sentences about Epoch prices are citations.

\subsubsection{3.4 Illustrative episodes, not
data}\label{illustrative-episodes-not-data}

Support desk. Completeness: the customer's issue is resolved, the reply
is sent, and the ticket state is Done. Further drafts do not change
\texttt{C}. A token counter can still increment. The episode is a
constructed example to show the variables. It is not a sample from a
help desk.

Care handoff. Completeness: critical fields filled, medication list
reconciled, receiving party confirmed, person in a safe state. Session
length is a different variable. The episode is constructed. No hospital
dataset was analyzed.

Coding chore. Completeness: the artifact matches a written oracle
(tests, or a bit-exact check). Regenerating a passing artifact does not
raise \texttt{C}.

\subsubsection{3.6 Fidelity criteria and stop
rules}\label{fidelity-criteria-and-stop-rules}

Shannon (1959) defines a rate-distortion function: the least description
rate that holds expected distortion at or below a chosen fidelity. The
object is communication under a constraint the user states. It is not a
demand curve for intelligence. The resemblance to this paper is only
structural. A completeness predicate \texttt{C} is a hard fidelity:
either the stated clauses hold, or they do not. A language-model loss is
closer to an average distortion on a training distribution. Falling
loss, and falling price at a fixed benchmark, can be read as movement
along a rate-resource curve of the kind Kaplan et al.~(2020) and
Hoffmann et al.~(2022) estimate for training. That movement does not
choose \texttt{C} for a particular chore. Confusing benchmark distortion
with chore completeness is the category error behind unbounded-seat
plans for closed work.

Cover and Thomas (2006) is the textbook path into rate-distortion if a
reader wants the formal development. It uses the existence of the
fidelity concept to keep two numbers apart: a score that can always be
improved by a fraction, and a predicate that is already true.

\subsubsection{3.7 Rebound taxonomy, without a new
estimate}\label{rebound-taxonomy-without-a-new-estimate}

The rebound canon is the frame for H2. Khazzoom (1980) showed that
mandated efficiency can raise use of the service, Brookes (1990) argued
the same at the level of the economy, and Saunders (1992) derived the
conditions in neoclassical growth. Sorrell and Dimitropoulos (2008) gave
the microeconomic definitions and their limits. Brockway, Sorrell,
Semieniuk, Heun and Court (2021) reviewed the evidence and found that
economy-wide rebound may erode more than half of the expected energy
savings. Task satiation does not contradict that finding. It is a
per-task stop, and the aggregate rebound of H2 is where the canon's
evidence applies. Lighting is the historical case to test it against:
Fouquet and Pearson (2006) trace seven centuries of the price and use of
light in the United Kingdom, and Tsao and colleagues (2010) find that
use of light rose with efficiency, a satiable service with large
rebound.

The rebound survey of Gillingham, Rapson, and Wagner separates channels
that a single slogan collapses. Direct rebound is additional use of the
same service when its effective cost falls. Indirect rebound is spending
of saved resources on other services. Economy-wide rebound runs through
prices and growth. The survey's policy domain is energy efficiency, not
tokens. The taxonomy still transfers as a warning about identification.
A measured fall in \texttt{p\_t} (claim S-2) identifies none of the
three channels by itself. Direct rebound on a chore with \texttt{C}
already true requires that buyers still pay for calls that do not change
\texttt{C}. Indirect and economy-wide channels can be large when new
tasks appear (Acemoglu and Restrepo, 2019). A paper that reported one
elasticity and called it ``the rebound'' would be under-specified. This
paper reports no elasticity.

The pattern is old. Jevons argued in 1865 that more efficient steam
engines would raise, not lower, Britain's total coal use. Alcott traces
that argument and the debate that followed (\emph{Ecological Economics}
54(1), 2005). Nordhaus measured the price of computation over more than
a century and found it falling by many orders of magnitude while total
computation rose (\emph{Journal of Economic History} 67(1), 2007). Agent
workloads already show the use side. Anthropic's engineering team
reported, for its own system, that agents use about 4 times the tokens
of a chat interaction and multi-agent systems about 15 times (13 June
2025). That figure is operator-reported, not audited. It is consistent
with per-task appetite rising as tasks change shape, which is exactly
what the split account in Section 3.3 must separate from post-completion
waste.

\subsubsection{3.8 Scarce joules and the economic stop (2025 to 2026
peers)}\label{scarce-joules-and-the-economic-stop-2025-to-2026-peers}

Epoch's price series (S-2) show synthesis getting cheaper in money.
Recent energy publications show joules remaining scarce, and often
growing in the aggregate.

Chung, Wu, Ma, and Chowdhury
(arXiv:\href{https://arxiv.org/abs/2601.22076}{2601.22076}, 2026)
diagnose inference energy across 1,858 configurations: task type alone
can move LLM energy per response by about 25×; video generation can
exceed image energy by two orders of magnitude. The \emph{Joule} article
``Energy use of AI inference, efficiency pathways, and test-time
scaling''
(\href{https://www.cell.com/joule/fulltext/S2542-4351(26)00114-5}{cell.com/joule/fulltext/S2542-4351(26)00114-5},
online 22 Apr 2026) estimates optimized frontier-scale queries in the
sub-watt-hour median range while reporting that long reasoning and
agentic queries can raise energy by more than an order of magnitude;
efficiency gains of 8 to 20× are line-of-sight, not free joules.
Luccioni, Strubell, and Crawford (FAccT 2025)
(\href{https://doi.org/10.1145/3715275.3732007\%5D(https://doi.org/10.1145/3715275.3732007)}{{[}doi:10.1145/3715275.3732007{]}(https://doi.org/10.1145/3715275.3732007))})
warn that efficiency without a stop can amplify total spend (Jevons).
Epoch AI's energy overview
(\href{https://epoch.ai/publications/ai-energy}{epoch.ai/publications/ai-energy},
updated Jul 2026) separates modest per-query costs from rapidly doubling
aggregate data-center power.

\textbf{Reading for this paper.} Free at the margin is a money-price
statement. Scarce joules are a physical statement. The operational stop
is task completion: when \texttt{C(z)\ =\ 1}, further synthesis does not
raise the written objective, and a stop on budget or policy grounds is
success. Estimates are not measured energy.

\subsubsection{3.9 Protocol for a done-detector study (not
executed)}\label{protocol-for-a-done-detector-study-not-executed}

The following protocol is a method for a future measurement. No step has
been run. It is written so the missing result has an exit criterion
rather than a slogan.

Population. Episodes drawn from one product surface, not a mix of
open-ended research and ticket closure. Each episode has a written
\texttt{C} with finitely many clauses, each clause observable by a
person who did not operate the model. Target size is a design choice;
this paper does not pretend a power calculation was done. A first usable
table needs enough positives and negatives that a rate is not a single
anecdote. The archive does not contain that table.

Label. Two annotators mark each clause true or false at the time the
system first declares done, and again at the end of the session.
Disagreements are adjudicated and counted. The detector's decision is
recorded separately from the label.

Primary rates.

\begin{itemize}
\tightlist
\item
  False done: system declares \texttt{C\ =\ 1} while any critical clause
  is false.
\item
  False continue: all critical clauses are true and the system issues
  further model or tool calls.
\item
  Session length and call count are covariates, not endpoints.
\end{itemize}

Secondary split, only if an energy account exists. Sum energy,
analytical or board-measured, on calls before the adjudicated
\texttt{t*} and after it. State the measurement class in the same
sentence as the number. An analytical OpCounter total is not a
substitute for a meter and is not comparable to Horowitz's tables.

Exclusion. Episodes with no written \texttt{C} are not labeled
``unsatiated.'' They are out of scope. Adversarial episodes are out of
scope for the satiation endpoint and may still be scored for physical
refuse if a certificate is present.

Failure of the protocol. If annotators cannot agree on clauses, the
predicate was not written. If the system has no declare-done event,
false continue is not identified. Publishing a model leaderboard instead
of the two rates would not execute this protocol.

\subsubsection{3.10 A numerical illustration of a bliss point (model,
not a
fit)}\label{a-numerical-illustration-of-a-bliss-point-model-not-a-fit}

Let one coordinate \texttt{x} be the count of model calls on a single
chore, and set

\begin{Shaded}
\begin{Highlighting}[]
\NormalTok{U(x) = {-}(x {-} 4)\^{}2}
\end{Highlighting}
\end{Shaded}

for a chore whose constructed maximum is at four calls. Then
\texttt{U(0)\ =\ -16}, \texttt{U(4)\ =\ 0}, \texttt{U(8)\ =\ -16}.
Marginal utility is \texttt{dU/dx\ =\ -2(x\ -\ 4)}, which is positive
before 4, zero at 4, and negative after. The illustration uses integers
so the sign change is visible without a fit. No dataset was used to
choose 4. Replacing 4 with another integer moves the peak and does not
change the definition. A product metric equal to \texttt{x} ranks
\texttt{x\ =\ 8} above \texttt{x\ =\ 4}. A product metric equal to
\texttt{U} does the reverse. Claim S-3 is this non-identity, here in a
fully specified toy function, and in Section 3.4 in constructed
episodes. Neither is a welfare estimate for a firm.

Token price can be attached without confusing units. Suppose a call
costs \texttt{p} in money and the buyer pays \texttt{p\ *\ x} while
value follows \texttt{U}. The buyer's surplus proxy
\texttt{U(x)\ -\ p\ x} still has a finite maximizer for any
\texttt{p\ \textgreater{}=\ 0}. A fall in \texttt{p} moves the
maximizer, which is the rebound logic on this toy, and the maximizer
remains finite. The toy does not say where the maximizer sits for any
real buyer. It blocks the inference that a lower \texttt{p} removes the
maximum.

Figure 1 draws the toy and the separate aggregate claim.

\includegraphics[width=0.95\linewidth,height=\textheight,keepaspectratio]{fig/satiation-rebound.png}

\emph{Figure 1. Task satiation and aggregate rebound. Model illustration
of the toy in this section, not a fit: the maximizer of \(U(x) - p\,x\)
is \(x = 4 - p/2\), finite at every price. The right panel is a
schematic; this paper estimates no rebound. A
\href{https://research.openie.dev/figures/satiation/satiation-rebound.png}{PNG
copy} is provided.}

\subsubsection{3.11 Reason codes as a partition, not a welfare
function}\label{reason-codes-as-a-partition-not-a-welfare-function}

Let \texttt{R} be the set of refuse codes on a trace. Partition
\texttt{R} into \texttt{R\_econ\ =\ \{policy,\ budget\_exceeded\}} and
\texttt{R\_phys\ =\ \{lut\_veto,\ energy\_veto,\ cbf\_veto\}}, matching
the schema's taxonomy. Other codes, if added later, have to be assigned
in writing before they are counted. A trace interval contributes to
post-done economic refuse only if \texttt{C} is already 1 and the code
is in \texttt{R\_econ}. It contributes to physical refuse if the code is
in \texttt{R\_phys}, whether or not \texttt{C} is 1. An unsafe action
after done is in both stories and should be counted in both columns, not
averaged. The MCP demo's irreversible refuse is a physical example:
energy and barrier codes, executor calls 0. It is not an
\texttt{R\_econ} example. A future episode with \texttt{policy} after a
merged patch would be the economic example. The repository does not
include that episode. The partition is how it would be scored.

\subsection{4. Evidence}\label{evidence}

\subsubsection{4.1 Claim S-1. Satiation is a defined economic
object}\label{claim-s-1.-satiation-is-a-defined-economic-object}

Andersen (2001), https://doi.org/10.1007/PL00003852, and the textbook
bliss construction in Section 2.1 support S-1. The claim is theoretical.
It is not a measured utility function for any user population in this
study.

\subsubsection{4.2 Claim S-2. Published inference prices fell sharply at
fixed
performance}\label{claim-s-2.-published-inference-prices-fell-sharply-at-fixed-performance}

Supported by Epoch AI (2025) and Emberson and Roodman (2026), with the
magnitudes those notes state (Section 2.5). The claim is about price
indices those authors publish. It is not a claim that every
organization's invoice fell by 13 times, and it is not a claim about
watts.

\subsubsection{4.3 Claim S-3. Engagement variables are not the
completeness
predicate}\label{claim-s-3.-engagement-variables-are-not-the-completeness-predicate}

Supported by the definitions in Section 3.1 and by the attention
literature's description of time-on-task objectives (Wu; Center for
Humane Technology overview). No causal estimate of welfare loss is
claimed. The logical claim is non-identity of the variables: token count
can increase while \texttt{C} stays constant. The constructed episodes
in Section 3.4 are illustrations of non-identity, not datasets.

\subsubsection{4.4 Claim S-4. Bicycle-mania prices detached from a
simple fundamental
story}\label{claim-s-4.-bicycle-mania-prices-detached-from-a-simple-fundamental-story}

Supported by Quinn (2019): rises over 200 percent and subsequent falls
over 75 percent in cycle shares, 1895 to 1900, not fully explained by
fundamentals or by a simple risk shift in that paper's tests. The claim
stops there. It does not quantify household satiation for bicycles.

\subsubsection{4.5 Claim S-5. U.S. smartphone ownership is high in the
cited fact
sheet}\label{claim-s-5.-u.s.-smartphone-ownership-is-high-in-the-cited-fact-sheet}

The project archive's reading of Pew's mobile fact sheet is ownership
near 91 percent for the 2024 and 2025 rows. The claim used in this paper
is diffusion at that order, conditional on the live sheet matching the
archive. It is a saturation fact about ownership, not a bliss-point
estimate.

\subsubsection{4.6 Claim S-6. Care is not priced as a pure software good
by the Baumol
framework}\label{claim-s-6.-care-is-not-priced-as-a-pure-software-good-by-the-baumol-framework}

Baumol (1967) supplies the framework. Acemoglu and Restrepo (2019)
supply the task language. The claim is that these frameworks do not
imply an end to care's labor, trust, or liability requirements when
token prices fall. The claim is not a forecast of health-care
expenditure.

\subsubsection{4.7 Claim S-7. A sell-side note questions spend versus
benefit}\label{claim-s-7.-a-sell-side-note-questions-spend-versus-benefit}

The Goldman Sachs page named in Section 2.6 exists and frames the
question in its title.

\subsubsection{4.9 Claim S-9. Physical law is not priced to zero by
Epoch}\label{claim-s-9.-physical-law-is-not-priced-to-zero-by-epoch}

Landauer (1961) and Horowitz (2014) support a permanent gap between an
ideal erasure bound and CMOS practice. S-9 says only that a citation of
Epoch does not discharge a citation of Landauer.

\subsubsection{4.10 Design rules, marked as
design}\label{design-rules-marked-as-design}

These rules follow if one accepts the definitions and wants a product
that optimizes \texttt{C} rather than tokens. They are not effects.

\begin{enumerate}
\def\labelenumi{\arabic{enumi}.}
\tightlist
\item
  Write \texttt{C} before adding model capacity.
\item
  Treat a synthesis price premium as temporary when Epoch-type series
  and open weights apply to that chore.
\item
  Do not describe modeled or simulated energy as measured energy.
\item
  Prefer a detector of \texttt{C} with an error table over an unlabeled
  increase in model size, once the chore is the product.
\item
  Require a separate physical safety check before any irreversible
  action.
\item
  Report pre-done and post-done energy separately when the split exists.
  It does not exist yet.
\item
  Do not apply the satiation predicate to adversarial or military races
  (Section 5.3).
\item
  If the product's objective is attention, say so. The completeness
  claims then do not apply.
\end{enumerate}

Kill tests for a roadmap, also design: a plan that assumes unbounded
appetite for a chore with a short \texttt{C}; a price that requires
permanent synthesis scarcity and no physical or liability wedge; a
safety story that uses only a confidence string; a ``free'' claim that
deletes joules; an agent loop with no \texttt{C}.

\subsubsection{4.11 What would be a result and is
not}\label{what-would-be-a-result-and-is-not}

A table of false-done and false-continue rates. A time series of calls
after \texttt{t*}. A meter reading of energy on a stated workload. A
structural rebound estimate. An AGI price premium estimated on a
satiated chore. None of these appear in Section 4's supported claims.

\subsubsection{4.12 Non-identity, stated as a comparison of
objectives}\label{non-identity-stated-as-a-comparison-of-objectives}

Let \texttt{x} be call count, \texttt{s} session length, \texttt{J\_pre}
and \texttt{J\_post} energy before and after adjudicated done, and
\texttt{C} the predicate. The following implications are definitional.

{\def\LTcaptype{none} % do not increment counter
\begin{longtable}[]{@{}
  >{\raggedright\arraybackslash}p{(\linewidth - 2\tabcolsep) * \real{0.4483}}
  >{\raggedright\arraybackslash}p{(\linewidth - 2\tabcolsep) * \real{0.5517}}@{}}
\toprule\noalign{}
\begin{minipage}[b]{\linewidth}\raggedright
Observation
\end{minipage} & \begin{minipage}[b]{\linewidth}\raggedright
Does not imply
\end{minipage} \\
\midrule\noalign{}
\endhead
\bottomrule\noalign{}
\endlastfoot
\texttt{x} increased & \texttt{C} increased \\
\texttt{s} increased & welfare increased \\
\texttt{p} fell (S-2) & \texttt{J\_post\ =\ 0} \\
\texttt{C\ =\ 1} & the detector's error rate is known \\
ownership of phones is high (S-5) & marginal value of another model call
is zero \\
cycle-share prices boomed and busted (S-4) & households reached a
transport bliss point \\
\end{longtable}
}

The right-hand column is the list of inferences this paper refuses. Each
left-hand observation is either a citation in Section 4 or a variable in
Section 3. The table is the discussion compressed into a check. A
sentence that asserts a right-hand claim has left the evidence.

\subsection{5. Discussion}\label{discussion}

\subsubsection{5.1 Rebound: the central identification
problem}\label{rebound-the-central-identification-problem}

Three results limit the natural experiment. Gundlach and colleagues
(2025) find that prices at fixed benchmark performance fall fast while
frontier prices do not, and Du (2026) finds a structural break and a
reasoning premium in price evolution, so price cuts arrive bundled with
model changes. Anthropic's Economic Index (September 2025) reports that
across tasks each 1\% increase in cost is associated with a 0.29\%
decrease in usage frequency, an inelastic cross-sectional figure, not a
price change over time. A provider-level design cannot identify
market-level rebound; the design needs market-level prices and
quantities.

Gillingham, Rapson, and Wagner survey rebound from energy-efficiency
policy. The mechanism is residual demand elasticity. If the money price
of synthesis falls, use rises where buyers still value another unit. At
a bliss point for a closed chore, the direct derivative of value with
respect to more synthesis on that chore is not positive, so direct
rebound on that chore is not implied. Indirect rebound, new tasks, and
attention shifting can remain. Acemoglu and Restrepo's new-task margin
is one place quantity can reappear. Claim S-2 therefore does not imply a
global cap on compute. It implies a cap only relative to a stated
\texttt{C}. Denying global rebound would require the split in Section
3.3 and a demand system. This paper does not deny global rebound. It
refuses to treat ``Jevons'' as a reason to avoid writing \texttt{C}.

Two hypotheses make the question testable.

\begin{itemize}
\tightlist
\item
  \textbf{H1. Post-completion spend.} The fraction of agent compute
  spent after \texttt{C(z)\ =\ 1}, joules or calls with
  \texttt{t\ \textgreater{}\ t*} over all joules or calls in the
  episode, is measurably above zero. H1 is the direct waste the
  completeness stop removes. It is estimable from agent logs with a
  labeled \texttt{t*} per episode (Section 3.3, item 1).
\item
  \textbf{H2. Jevons in tasks.} H2 is the direct-rebound elasticity of
  Khazzoom (1980), stated for tasks. When the price of synthesis falls,
  the elasticity of the number of tasks run exceeds the fall in spend
  per task. Then total spend rises even though each task is satiated. H2
  is aggregate rebound.
\end{itemize}

\textbf{Identification.} H1 needs no price variation: a labeled corpus,
the detector error rates of Section 3.3 item 2, and the split account.
H2 needs exogenous price changes. Provider price cuts at fixed model
quality, announced on a known date, are natural experiments. The design
is a panel of accounts with price paid, tasks opened, calls per task and
a completeness label per episode, before and after each cut. Comparing
accounts on a repriced model with accounts on models whose price did not
move gives a difference-in-differences estimate of the task-count
elasticity and the per-task spend response separately. Spend per task
can fall while task count rises. Both numbers are needed. The
completeness stop is a claim about the first. Jevons is a claim about
the second. Neither was estimated here.

\subsubsection{5.2 Status and creative
goods}\label{status-and-creative-goods}

Some goods are valued because they are not finished: status, open-ended
art, research whose \texttt{C} is not a finite checklist. The definition
in Section 3.1 does not force a bliss point onto those goods. It forces
a segmentation. A finite chore and an open-ended practice can be sold by
the same firm and should not share a demand assumption. Pricing a finite
chore as if it were an open-ended capability race is a category error,
not a finding about AGI.

\subsubsection{5.3 Adversarial demand}\label{adversarial-demand}

If the other party's gain is this party's loss, ``enough'' is not
defined by a local checklist. Arms races and some security contests are
in that class. Physical commit constraints can still apply: an unsafe
action can be refused while the strategic demand remains open. Scope is
part of the claim, not a footnote that cancels it.

\subsubsection{5.4 Incomplete knowledge of
done}\label{incomplete-knowledge-of-done}

People misjudge completion. Detectors will false-done and
false-continue. That is an argument for measuring detector error, which
Section 3.3 lists as missing. It is not an argument that session length
is a substitute for \texttt{C} in a domain where a missed handoff has a
liable party. Where no predicate can be written, this paper's
operational claim is silent.

\subsubsection{5.5 Baumol is not solved}\label{baumol-is-not-solved}

A model that drafts a note faster can raise throughput of notes and
leave examination, judgment, and responsibility with a person.
Expenditure shares can still move as Baumol described if the human
residual does not speed up. Saying otherwise would be a forecast this
bibliography does not contain. The permitted statement is S-6.

\subsubsection{5.6 Free synthesis is not free
actuation}\label{free-synthesis-is-not-free-actuation}

A falling money price for inference says nothing about the energy of
moving a physical system. Moving a mass \(m\) through a speed change
\(\Delta v\) from rest takes at least \(\tfrac12 m (\Delta v)^2\) joules
of work, whatever the price of the instruction that commanded it. Epoch
prices sit on a different axis from that work. Collapsing them into
``free AI'' deletes the measurement class.

\subsubsection{5.7 Confidence is not a
certificate}\label{confidence-is-not-a-certificate}

A high model score can order proposals. It is not a completeness
predicate: a score needs a threshold and an error table (false and true
positive rates against a labeled set) before it decides \texttt{C}.

\subsubsection{5.8 Sell-side skepticism is not a
measurement}\label{sell-side-skepticism-is-not-a-measurement}

S-7 prevents a literature review from ignoring a major public doubt
about spend and benefit. It does not let a research note replace a meter
or a demand estimate. If the note's figures are quoted later, the
primary PDF has to be the source, and the label is sell-side estimate.

\subsubsection{5.9 What a structural demand estimate would
require}\label{what-a-structural-demand-estimate-would-require}

S-2 is a supply-price fact from secondary published series. A demand
estimate is a different regression. At minimum it needs a quantity
(calls, tasks closed, or both), a price, and a way to separate movement
along a demand curve from movement of the curve. Epoch's benchmark price
is not a firm-level invoice, and a firm's invoice is not \texttt{C}.
Linking them requires a panel this study does not have: accounts that
record price paid, calls made, and an adjudicated completeness label on
the same episode. Without that panel, the slope of demand after
\texttt{t*} is not identified. The toy surplus in Section 3.10 shows
only that a quadratic bliss point keeps a finite optimum as \texttt{p}
falls. It is not the missing regression.

Segmentation belongs in the same design. Code each product as
finite-\texttt{C} or open-ended before pooling. Pooling a ticket desk
with an open research tool averages two demand regimes and can
manufacture an appearance of unbounded appetite. The military exclusion
in Section 5.3 is the same rule applied to a domain where the opponent's
objective prevents a local bliss point. A table with those codes, even
with no prices yet, would be more informative than a single
total-addressable-market number. The table was not built.

\subsubsection{\texorpdfstring{5.10 Why no multiple of
\texttt{kT\ ln\ 2} is
stated}{5.10 Why no multiple of kT ln 2 is stated}}\label{why-no-multiple-of-kt-ln-2-is-stated}

Landauer's result is an inequality for an idealized erasure at
temperature \texttt{T}. Horowitz (2014) reports energies for particular
CMOS operations and for memory traffic, in picojoules, under the
assumptions of that tutorial. Dividing one by the other produces a
number only after someone chooses the operation, the activity factor,
the temperature, and the device. None of those choices is measured in
this paper. Stating a fold above the Landauer bound would pretend those
choices had been made. The scientific statement is the inequality
direction: practical irreversible CMOS in Horowitz's account sits far
above the bound, and a price index does not move a device onto the
bound.

\subsection{6. Limits and threats to
validity}\label{limits-and-threats-to-validity}

Construct validity. Operational satiation depends on \texttt{C}. A
hostile or vague \texttt{C} makes the predicate arbitrary. The paper's
defense is to require \texttt{C} to be written in observable clauses.
Many real products do not have that writing yet. For those products the
theory does not yet apply empirically.

Internal validity of the price claim. S-2 inherits whatever
identification Epoch used (quality adjustment, benchmark choice, posted
prices versus realized invoices). This paper did not audit the
microdata. A revision of Epoch's series would revise S-2's magnitudes
and might not revise the direction. S-2 should be re-read against the
primary pages before it is used as an executive figure.

External validity. Work and care, as defined, exclude military rivalry
and may exclude research. Phone ownership does not transfer to agent
products. Bicycle share prices do not transfer to software margins. The
constructed episodes are not a sample.

Measurement validity. No split joules, no detector confusion matrix, and
no energy meter on any device. Vendor hardware specifications are not
energy measurements.

Omitted variable. New tasks can absorb spending released by cheap
synthesis. Section 5.1 states this. A roadmap that treats per-chore
satiation as a forecast of aggregate compute demand is not licensed by
this paper.

Publication bias and document type. Wu is a book. The Humane Technology
page is advocacy. Goldman Sachs is sell-side. Quinn and Epoch and Baumol
and Acemoglu and Restrepo and Landauer and Horowitz are the load-bearing
citations, and they support different sentences. The paper's structure
is an attempt to keep those sentences from being averaged into a mood.

Falsifiers. S-1 fails if the cited theory does not contain a zero
marginal-utility region. It does. S-2 fails if the primary Epoch pages
do not report declines of that kind. A future labeled corpus in which
extra calls after a correct \texttt{C} systematically raise an
independently specified value would falsify satiation for that corpus.
One such corpus would not falsify the definition. It would bound its
domain.

Vendor hardware specifications list capacities, not switching activity,
static power or the energy of a workload. An energy report has to name
the hardware, the clock, the workload, the meter and the duration.

\subsection{7. Conclusion}\label{conclusion}

Satiation is a standard object in the theory cited, and it has a narrow
operational form: a written completeness predicate that extra synthesis
does not advance. Epoch's published series document a sharp fall in the
money price of fixed-performance inference. That fall supports a
free-at-the-margin description of digital chores and does not support
free joules, free plant motion, or the end of care's cost structure.
Bicycle share prices, smartphone diffusion, and attention markets are
weaker and stronger in the specific ways Section 4 states. The missing
scientific objects are a labeled done detector, a pre-done versus
post-done energy split, and a demand estimate that can see rebound.
Until they exist, the paper's empirical content is the cited record, and
its formal content is the predicate.

\subsection{References}\label{references}

Cover, T. M., and Thomas, J. A. Elements of Information Theory, 2nd
ed.~Wiley, 2006. https://doi.org/10.1002/047174882X

Hoffmann, J., et al.~Training compute-optimal large language models.
arXiv:2203.15556.

Kaplan, J., et al.~Scaling laws for neural language models.
arXiv:2001.08361.

Shannon, C. E. Coding theorems for a discrete source with a fidelity
criterion. IRE National Convention Record, 1959.

Acemoglu, D., and Restrepo, P. Automation and new tasks: how technology
displaces and reinstates labor. Journal of Economic Perspectives 33(2),
2019. https://www.aeaweb.org/articles?id=10.1257/jep.33.2.3

Ames, A. D., et al.~Control barrier functions: theory and applications.
ECC 2019. https://doi.org/10.23919/ECC.2019.8796030

Andersen, E. S. Satiation in an evolutionary model of structural
economic dynamics. Journal of Evolutionary Economics, 2001.
https://doi.org/10.1007/PL00003852

Baumol, W. J. Macroeconomics of unbalanced growth: the anatomy of urban
crisis. American Economic Review 57, 415-426, 1967.

Bennett, C. H. Logical reversibility of computation. IBM Journal of
Research and Development, 1973. https://doi.org/10.1147/rd.176.0525

Center for Humane Technology. The attention economy.
https://www.humanetech.com/youth/the-attention-economy

Emberson, L., and Roodman, D. / Epoch AI. The plunging price of thought.
22 September 2026.
https://epoch.ai/publications/the-plunging-price-of-thought

Epoch AI. LLM inference price trends. 12 March 2025.
https://epoch.ai/data-insights/llm-inference-price-trends

Gillingham, K., Rapson, D., and Wagner, G. The rebound effect and energy
efficiency policy. Review of Environmental Economics and Policy 10(1),
68-88, 2016. https://doi.org/10.1093/reep/rev017 Working text:
https://gwagner.com/wp-content/uploads/Gillingham-Rapson-Wagner-2015-Rebound-Effect.pdf

Alcott, B. Jevons' paradox. Ecological Economics 54(1), 9-21, 2005.
https://doi.org/10.1016/j.ecolecon.2005.03.020

Anthropic. How we built our multi-agent research system. 13 June 2025.
https://www.anthropic.com/engineering/multi-agent-research-system

Jevons, W. S. The Coal Question. London: Macmillan, 1865.

Keynes, J. M. Economic possibilities for our grandchildren. 1930.
Reprinted in Essays in Persuasion, 1931.

Nordhaus, W. D. Two centuries of productivity growth in computing.
Journal of Economic History 67(1), 128-159, 2007.
https://doi.org/10.1017/S0022050707000058

Goldman Sachs. Gen AI: too much spend, too little benefit?
https://www.goldmansachs.com/insights/top-of-mind/gen-ai-too-much-spend-too-little-benefit

Horowitz, M. Computing's energy problem (and what we can do about it).
ISSCC 2014. https://doi.org/10.1109/ISSCC.2014.6757323

Landauer, R. Irreversibility and heat generation in the computing
process. IBM Journal of Research and Development, 1961.
https://doi.org/10.1147/rd.53.0183

Laya. System One models and typed decision engines (comparison table,
checked 2 October 2026). https://laya-ai.com/system-one-models

Convai Innovations. Laya model card.
https://huggingface.co/convaiinnovations/laya

TypeSafe AI. Introducing System One Models and Jev. 15 September 2026.
https://typesafe.ai/blog/introducing-system-one-models-and-jev

Ma, S., et al.~The era of 1-bit LLMs. arXiv:2402.17764.

Pew Research Center. Mobile fact sheet.
https://www.pewresearch.org/internet/fact-sheet/mobile/

Quinn, W. Technological revolutions and speculative finance: evidence
from the British Bicycle Mania. Cambridge Journal of Economics 43(2),
271-294, 2019. https://doi.org/10.1093/cje/bey029

Wei, J., et al.~T-MAC. arXiv:2407.00088.

Wu, T. The Attention Merchants. Knopf.
https://www.penguinrandomhouse.com/books/234876/the-attention-merchants-by-tim-wu/

Chung, J.-W., Wu, R., Ma, J. J., and Chowdhury, M. Where Do the Joules
Go? Diagnosing Inference Energy Consumption. arXiv:2601.22076.
https://arxiv.org/abs/2601.22076

Energy use of AI inference, efficiency pathways, and test-time scaling.
\emph{Joule}, 2026.
https://www.cell.com/joule/fulltext/S2542-4351(26)00114-5

Luccioni, A. S., Strubell, E., and Crawford, K. From efficiency gains to
rebound effects: the problem of Jevons' Paradox in AI's polarized
environmental debate. \emph{FAccT} 2025.
https://doi.org/10.1145/3715275.3732007

Ostrovsky, N. / Epoch AI. What you need to know about AI energy use.
Updated 14 July 2026. https://epoch.ai/publications/ai-energy

Anthropic. Anthropic Economic Index report: uneven geographic and
enterprise AI adoption. September 2025.
https://www.anthropic.com/research/anthropic-economic-index-september-2025-report

Bai, L., Huang, Z., Wang, X., et al.~How do AI agents spend your money?
Analyzing and predicting token consumption in agentic coding.
arXiv:2604.22750, 2026. https://arxiv.org/abs/2604.22750

Brockway, P. E., Sorrell, S., Semieniuk, G., Heun, M. K., and Court, V.
Energy efficiency and economy-wide rebound effects: a review of the
evidence and its implications. Renewable and Sustainable Energy Reviews
141, 110781, 2021. https://doi.org/10.1016/j.rser.2021.110781

Brookes, L. The greenhouse effect: the fallacies in the energy
efficiency solution. Energy Policy 18(2), 199-201, 1990.
https://doi.org/10.1016/0301-4215(90)90145-T

Chen, X., Xu, J., Liang, T., et al.~Do NOT think that much for 2+3=? On
the overthinking of o1-like LLMs. arXiv:2412.21187, 2024.
https://arxiv.org/abs/2412.21187

Cuadron, A., Li, D., Ma, W., et al.~The danger of overthinking:
examining the reasoning-action dilemma in agentic tasks.
arXiv:2502.08235, 2025. https://arxiv.org/abs/2502.08235

Du, M. Tiered super-Moore's law: price evolution, production frontiers,
and market competition in large language models. arXiv:2603.28576, 2026.
https://arxiv.org/abs/2603.28576

Fouquet, R., and Pearson, P. J. G. Seven centuries of energy services:
the price and use of light in the United Kingdom (1300-2000). The Energy
Journal 27(1), 2006.
https://doi.org/10.5547/ISSN0195-6574-EJ-Vol27-No1-8

Gundlach, H., Lynch, J., Mertens, M., et al.~The price of progress:
price performance and the future of AI. arXiv:2511.23455, 2025.
https://arxiv.org/abs/2511.23455

Khazzoom, J. D. Economic implications of mandated efficiency in
standards for household appliances. The Energy Journal 1(4), 1980.
https://doi.org/10.5547/ISSN0195-6574-EJ-Vol1-No4-2

Masanet, E., Shehabi, A., Lei, N., Smith, S., and Koomey, J.
Recalibrating global data center energy-use estimates. Science
367(6481), 984-986, 2020. https://doi.org/10.1126/science.aba3758

Saunders, H. D. The Khazzoom-Brookes postulate and neoclassical growth.
The Energy Journal 13(4), 1992.
https://doi.org/10.5547/ISSN0195-6574-EJ-Vol13-No4-7

Sorrell, S., and Dimitropoulos, J. The rebound effect: microeconomic
definitions, limitations and extensions. Ecological Economics 65(3),
636-649, 2008. https://doi.org/10.1016/j.ecolecon.2007.08.013

Tsao, J. Y., Saunders, H. D., Creighton, J. R., Coltrin, M. E., and
Simmons, J. A. Solid-state lighting: an energy-economics perspective.
Journal of Physics D: Applied Physics 43(35), 354001, 2010.
https://doi.org/10.1088/0022-3727/43/35/354001

Witt, U. Learning to consume: a theory of wants and the growth of
demand. Journal of Evolutionary Economics 11(1), 23-36, 2001.
https://doi.org/10.1007/PL00003851

\subsection{Appendix A. Decision rules and
non-results}\label{appendix-a.-decision-rules-and-non-results}

{\def\LTcaptype{none} % do not increment counter
\begin{longtable}[]{@{}
  >{\raggedright\arraybackslash}p{(\linewidth - 4\tabcolsep) * \real{0.1818}}
  >{\raggedright\arraybackslash}p{(\linewidth - 4\tabcolsep) * \real{0.5000}}
  >{\raggedright\arraybackslash}p{(\linewidth - 4\tabcolsep) * \real{0.3182}}@{}}
\toprule\noalign{}
\begin{minipage}[b]{\linewidth}\raggedright
Id
\end{minipage} & \begin{minipage}[b]{\linewidth}\raggedright
Statement
\end{minipage} & \begin{minipage}[b]{\linewidth}\raggedright
Class
\end{minipage} \\
\midrule\noalign{}
\endhead
\bottomrule\noalign{}
\endlastfoot
R1 & Write the completeness predicate before adding capacity. &
Design \\
R2 & Expect synthesis price premia to erode where Epoch-type evidence
applies. & Design, conditional on S-2 \\
R3 & Do not claim free energy or free actuation. & Constraint from
S-9 \\
R4 & Separate economic stop reasons from physical stop reasons. &
Design \\
R5 & Publish detector error rates before treating a stop as accurate. &
Not yet a result \\
R6 & Split energy before and after done. & Not yet a result \\
R7 & Keep adversarial races outside the satiation claim. & Scope \\
R8 & Re-read Pew and Epoch primaries before quoting a finer magnitude. &
Method \\
\end{longtable}
}

\subsection{Appendix C. Claim map}\label{appendix-c.-claim-map}

Supported now: S-1 through S-9, with the scope sentences in Section 4.

Not supported: an end to Baumol cost disease; free plant motion;
satiation of military demand; a numerical global rebound of zero;
Landauer-scale operation of a language model; board watts; an AGI price
as a fact about satiated chores; a universal distillation lag; silicon
performance leadership.

Open measurements: Section 3.3.

\end{document}
